% ============================================================================
% Fichier  : partie2/P02_C07_cobit_conception.tex
% Titre    : Concevoir un Système de Gouvernance sur Mesure
% Auteur   : MINKA MI NGUIDJOÏ Thierry Emmanuel
% Version  : 1.0 -- octobre 2026
% ============================================================================
% Source : COBIT 2019 Framework: Introduction and Methodology, chapitres 4.5
% et 7 (dossier Normes). Le COBIT 2019 Design Guide (outil de calcul
% officiel) n'est pas dans le dossier : la méthode de calcul de la
% section 4 est une méthode pédagogique simplifiée, signalée comme telle.
% Laboratoires associés : GCI-LAB-02 (facteurs de conception) et GCI-LAB-03
% (cascade des objectifs).
% ============================================================================

\chapter{Concevoir un Système de Gouvernance sur Mesure}
\label{chap:cobit-conception}
\chaptermark{Concevoir un système sur mesure}

\epigraph{%
    « Il n'existe pas de formule magique. La conception finale est une
    décision au cas par cas. »%
}{D'après ISACA, \textit{COBIT 2019 Framework: Introduction and
Methodology} (2018, notre traduction)}

\niveauChapitre{Intermédiaire}{%
    Maîtrise du Chapitre~\ref{chap:cobit-objectifs} (les 40~objectifs et la
    cascade des objectifs).%
}

\begin{boxobjectifs}
\begin{itemize}
    \item Appliquer les onze facteurs de conception.
    \item Dérouler le processus de conception en quatre étapes.
    \item Calculer un profil d'objectifs prioritaires et en fixer les
          niveaux de capacité cibles.
    \item Justifier les arbitrages devant une direction générale.
\end{itemize}
\end{boxobjectifs}

\bigskip

% ============================================================================
\section*{Question directrice : dix objectifs, mais jusqu'où~?}
% ============================================================================

Au Chapitre~\ref{chap:cobit-objectifs}, la BANQUE SANAGA a retenu dix
objectifs prioritaires. Le directeur général, M.~FOTSO, pose aussitôt la
question suivante~: « Faut-il porter les dix au même niveau d'exigence~?
Notre budget de mise en œuvre ne le permettra pas. » À Bafoussam, la
fondatrice de NOUN PAY, Mme~KAMGA, pose la question inverse~: « Avec
quarante-cinq salariés, avons-nous seulement besoin de COBIT~? »

Les deux questions ont la même réponse~: COBIT~2019 ne s'applique pas
tel quel, il se \textbf{conçoit}. Les quarante objectifs sont, dans le
modèle, d'importance égale~; ce sont les caractéristiques de chaque
organisation, les \emph{facteurs de conception}, qui les hiérarchisent.
Ce chapitre présente leurs effets (section~\ref{sec:c07-effets}), les
onze facteurs (section~\ref{sec:c07-facteurs}), la démarche en quatre
étapes (section~\ref{sec:c07-etapes}), une méthode de calcul
(section~\ref{sec:c07-calcul}) et l'art de l'arbitrage
(section~\ref{sec:c07-arbitrage}). Les laboratoires GCI-LAB-02 et GCI-LAB-03 du
NDAP SUKULU en sont l'application.

\separateur

% ============================================================================
\section{Les Trois Effets des Facteurs de Conception}
\label{sec:c07-effets}
% ============================================================================

Le modèle central de COBIT~2019 traite ses quarante objectifs comme
intrinsèquement équivalents~: il n'existe entre eux aucun ordre de priorité
naturel. Les facteurs de conception rompent cette équivalence de trois
manières~\cite{isaca2018cobitintro}.

\begin{enumerate}
    \item \textbf{La priorité des objectifs et leur capacité cible.} Certains
          objectifs deviennent plus importants que d'autres, et quelques-uns
          peuvent même devenir négligeables. En pratique, l'importance se
          traduit par un \emph{niveau de capacité cible} plus élevé. Une
          organisation très averse au risque, par exemple, donnera plus de
          poids à EDM03, APO12, APO13 et DSS05.
    \item \textbf{La variation des composantes.} Certains facteurs changent
          l'importance d'une composante, ou en exigent une variante. Une
          petite organisation n'a pas besoin de toutes les structures du
          modèle~; une organisation très réglementée accorde plus
          d'importance aux procédures écrites et à la fonction conformité.
    \item \textbf{Le recours à un domaine d'intérêt.} Certains facteurs
          (menaces, méthodes de développement, taille) appellent un contenu
          adapté, décrit dans un domaine d'intérêt~: DevOps, petites et
          moyennes entreprises, cybersécurité.
\end{enumerate}

\begin{boxdef}
\termecle{Niveau de capacité cible}\index{niveau de capacité cible}
(\textit{target capability level})~: niveau de capacité, de~0 à~5, que
l'organisation décide d'atteindre pour un processus donné. Il traduit
l'importance de l'objectif pour elle. Le Chapitre~\ref{chap:cobit-capacite}
définit chaque niveau et la manière de mesurer la capacité constatée.
\end{boxdef}

% ============================================================================
\section{Les Onze Facteurs de Conception}
\label{sec:c07-facteurs}
% ============================================================================

\subsection{Le tableau des facteurs}

COBIT~2019 retient onze facteurs, dont les huit derniers sont parfois
appelés \emph{facteurs contextuels}. Pour chacun, il propose des valeurs
types~\cite{isaca2018cobitintro}. Le tableau suivant les applique à deux
organisations du cours.

\begin{tableaucours}{L{2.7cm}L{3.5cm}L{3.0cm}L{3.0cm}}{%
    \tableauHeader{Facteur} &
    \tableauHeader{Valeurs types} &
    \tableauHeader{BANQUE SANAGA} &
    \tableauHeader{ANPP}
}
1. Stratégie & Croissance, innovation, maîtrise des coûts, service et
stabilité & Service et stabilité (après la panne) & Service et stabilité \\
2. Objectifs d'entreprise & EG01 à EG13 & EG06, EG02, EG03 & EG12, EG09,
EG03 \\
3. Profil de risque & 19~catégories de risques liés aux SI & Incidents
d'exploitation, fournisseurs, non-conformité & Projets, décisions
d'investissement, fournisseurs \\
4. Problèmes liés aux SI & 20~problèmes types & Incidents majeurs,
défaillance du prestataire, constats d'audit répétés & Projets en retard,
doublons, dépenses hors budget \\
5. Paysage des menaces & Normal, élevé & Élevé (secteur bancaire) &
Normal \\
6. Exigences de conformité & Faibles, normales, élevées & Élevées (COBAC,
loi n°~2010/012) & Élevées (gestion publique) \\
7. Rôle des SI & Soutien, usine, redressement, stratégique & Stratégique
& Usine \\
8. Approvisionnement & Externalisé, nuage, interne, hybride & Hybride &
Externalisé \\
9. Méthodes de mise en œuvre & Agile, DevOps, traditionnelle, hybride &
Traditionnelle & Traditionnelle \\
10. Adoption des technologies & Pionnier, suiveur, tardif & Suiveur &
Tardif \\
11. Taille & Grande (plus de 250), PME (50 à 250) & Grande & Grande \\
\end{tableaucours}
\vspace{-0.8em}
\begin{center}{\small\itshape Tableau~7.1. Les onze facteurs de conception
appliqués à deux organisations (valeurs types d'après COBIT~2019, notre
traduction)}\end{center}

\subsection{Quatre facteurs à bien comprendre}

\paragraph{Le rôle des SI.} COBIT reprend une typologie classique en
quatre rôles. Les SI ont un rôle de \emph{soutien} quand ils ne sont
cruciaux ni pour le fonctionnement ni pour l'innovation~; d'\emph{usine}
quand leur panne arrête immédiatement l'activité, sans qu'ils soient un
moteur d'innovation~; de \emph{redressement} quand ils sont vus comme un
moteur d'innovation sans que l'activité courante en dépende encore~;
\emph{stratégiques} quand ils sont critiques à la fois pour fonctionner et
pour innover. L'ANPP est une usine~: une panne arrête les paiements, mais
l'agence n'innove pas par le numérique. La BANQUE SANAGA est stratégique.

\paragraph{Les problèmes liés aux SI.} Le profil de risque regarde vers
l'avenir~; les problèmes regardent ce qui s'est déjà produit. COBIT en
propose une liste type~: frustrations entre métiers et DSI, incidents
majeurs, problèmes de service des prestataires, constats d'audit répétés,
dépenses informatiques cachées, projets en retard, solutions développées
par les métiers sans la DSI, méconnaissance des règles de protection de la
vie privée. C'est par ce facteur que la panne de la BANQUE SANAGA remonte
BAI06 et BAI07 dans ses priorités (Chapitre~\ref{chap:cobit-objectifs}).

\paragraph{Les exigences de conformité.} Elles mesurent le poids des
obligations légales, réglementaires et contractuelles. Au Cameroun, elles
sont élevées pour les banques, les opérateurs de communications
électroniques et, depuis 2026, pour toute organisation qui traite des
données personnelles à grande échelle.

\paragraph{La taille.} COBIT distingue les grandes organisations (plus de
250~salariés à temps plein) et les petites et moyennes (de 50 à~250).
Les très petites organisations, de moins de 50~salariés, ne sont pas
couvertes par le modèle.

\subsection{L'ancrage juridique camerounais}

\begin{boxjuris}
\textbf{\loicam{n°~2010/012} du 21~décembre~2010 relative à la
cybersécurité et la cybercriminalité au Cameroun, article~31, alinéa~1}
(version anglaise officielle)~:

\medskip
\textit{« Electronic communication networks and information systems
content providers shall be bound to ensure the availability of material,
as well as the data stored in their installations. »}

\medskip
Soit~: les fournisseurs de contenus des réseaux de communications
électroniques et des systèmes d'information sont tenus d'assurer la
disponibilité des contenus et des données stockées dans leurs
installations.

\medskip
\textbf{\loicam{n°~2024/017} du 23~décembre~2024 relative à la protection
des données à caractère personnel.} D'après les présentations publiées
lors de son adoption~\cite{cio2026donnees}, cette loi s'applique
pleinement depuis le 23~juin~2026, au terme d'un délai de mise en
conformité de dix-huit mois. Elle institue une Autorité de protection des
données à caractère personnel et encadre strictement le transfert des
données hors du territoire national.

\medskip
\textbf{Portée pour la conception.} L'obligation de disponibilité de
l'article~31, l'obligation d'audit des articles~13 et~32
(Chapitre~\ref{chap:prologue}) et, depuis 2026, les obligations relatives
aux données personnelles justifient de classer les \emph{exigences de
conformité} en « élevées » pour la BANQUE SANAGA et pour NOUN PAY. Ce
classement relève en priorité DSS04 (continuité), MEA03 (conformité aux
exigences externes) et APO14 (données). Le Chapitre~\ref{chap:droit}
étudiera ces textes en détail.
\end{boxjuris}
% A VERIFIER au Journal officiel : articles de la loi n° 2024/017 (date
% d'application, autorité, transferts) ; sources utilisées : presse
% spécialisée (CIO Mag, Droit Médias Finance).

% ============================================================================
\section{La Démarche en Quatre Étapes}
\label{sec:c07-etapes}
% ============================================================================

COBIT~2019 propose une démarche de conception en quatre étapes, chacune
décomposée en sous-étapes~\cite{isaca2018cobitintro}.

\noindent\begin{minipage}{\linewidth}
\begin{center}
\begin{tikzpicture}[
    etape/.style = {
        draw=CouleurPrimaire, fill=CouleurPrimaire, rounded corners=4pt,
        text width=2.7cm, minimum height=1.0cm, align=center,
        font=\footnotesize\bfseries\color{white}
    },
    detail/.style = {
        draw=CouleurBordDef, fill=CouleurDef, rounded corners=3pt,
        text width=2.7cm, align=left, anchor=north,
        font=\scriptsize\color{CouleurPrimaire}
    },
    fleche/.style = {-{Stealth[length=5pt]}, thick, color=CouleurSecondaire}
]
\node[etape] (e1) at (0,0)    {1. Comprendre le contexte et la stratégie};
\node[etape] (e2) at (3.4,0)  {2. Déterminer le périmètre initial};
\node[etape] (e3) at (6.8,0)  {3. Affiner le périmètre};
\node[etape] (e4) at (10.2,0) {4. Conclure la conception};
\draw[fleche] (e1) -- (e2);
\draw[fleche] (e2) -- (e3);
\draw[fleche] (e3) -- (e4);
\node[detail] at (0,-0.75)    {Stratégie, objectifs d'entreprise, profil de
    risque, problèmes liés aux SI};
\node[detail] at (3.4,-0.75)  {Appliquer ces quatre facteurs, dont la
    cascade des objectifs};
\node[detail] at (6.8,-0.75)  {Menaces, conformité, rôle des SI,
    approvisionnement, méthodes, adoption, taille};
\node[detail] at (10.2,-0.75) {Résoudre les conflits de priorités,
    arrêter la conception};
\end{tikzpicture}
\end{center}
\vspace{-0.4em}
\begin{center}{\small\itshape Figure~7.1. La démarche de conception de
COBIT~2019 (d'après ISACA, notre traduction)}\end{center}
\end{minipage}

\medskip

Les deux premières étapes travaillent sur les quatre premiers facteurs~:
la première les \emph{comprend}, la deuxième les \emph{applique} pour
dégager un premier périmètre. La troisième étape affine ce périmètre avec
les sept facteurs contextuels. La quatrième résout les contradictions et
arrête la conception.

COBIT recommande de reporter toutes les indications obtenues sur un
\termecle{canevas de conception}\index{canevas de conception}
(\textit{design canvas}), puis de résoudre, autant que possible, les
conflits qui y apparaissent. Ces conflits sont inévitables~: avec onze
facteurs et des correspondances génériques, certaines indications se
contredisent.

% ============================================================================
\section{Calculer un Profil d'Objectifs}
\label{sec:c07-calcul}
% ============================================================================

\subsection{Une méthode pédagogique}

Le guide de conception de COBIT~2019 (\textit{Design Guide}) et son outil
de calcul attribuent, pour chaque facteur et chaque valeur, des
coefficients d'importance à chacun des quarante objectifs. Ces tables ne
sont pas reproduites ici. Ce cours utilise une \textbf{méthode
pédagogique simplifiée}, qui suit la même logique~: chaque facteur
contribue à l'importance de chaque objectif, et l'importance se traduit en
niveau de capacité cible. Le laboratoire GCI-LAB-02 applique cette méthode~;
en entreprise, on utilisera l'outil officiel.

\begin{boxalert}
\textbf{À ne pas confondre.} Les poids, les notes et les seuils de la
méthode ci-dessous sont des conventions du cours, fixées pour
l'apprentissage. Ce ne sont pas les coefficients officiels du
\textit{Design Guide}. Un rapport professionnel qui les présenterait comme
tels commettrait une erreur de source.
\end{boxalert}

\subsection{Les quatre temps du calcul}

\begin{enumerate}
    \item \textbf{Retenir les facteurs déterminants} et leur donner un poids
          (la somme des poids vaut 100~\%). Le choix des poids est une
          décision de gouvernance.
    \item \textbf{Noter, pour chaque objectif, la contribution de chaque
          facteur}, de~0 (aucune) à~3 (très forte), en justifiant chaque
          note.
    \item \textbf{Calculer le score d'importance} de l'objectif~:
          \[
              S = \frac{\sum_{f} w_f \, n_f}{3} \times 100,
          \]
          où $w_f$ est le poids du facteur et $n_f$ la note. Le score va
          de~0 à~100.
    \item \textbf{Traduire le score en capacité cible} selon la convention
          du cours~: 75 et plus, niveau~4~; de 50 à~74, niveau~3~; de 25
          à~49, niveau~2~; moins de~25, niveau~1.
\end{enumerate}

\subsection{Application à la BANQUE SANAGA}

\begin{boxscenario}{Opération SANAGA. Le profil de la banque}
Le comité stratégique des SI retient cinq facteurs et leurs poids~:
cascade des objectifs (25~\%), profil de risque (25~\%), problèmes liés aux
SI (20~\%), exigences de conformité (15~\%), approvisionnement (15~\%).

\medskip
\begin{tableaucours}{L{1.5cm}cccccrc}{%
    \tableauHeader{Objectif} &
    \tableauHeader{Casc.} &
    \tableauHeader{Risq.} &
    \tableauHeader{Prob.} &
    \tableauHeader{Conf.} &
    \tableauHeader{Appr.} &
    \tableauHeader{Score} &
    \tableauHeader{Cible}
}
DSS04 & 3 & 3 & 3 & 3 & 1 & 90,0 & 4 \\
EDM03 & 3 & 3 & 2 & 3 & 1 & 83,3 & 4 \\
APO12 & 3 & 3 & 2 & 2 & 1 & 78,3 & 4 \\
BAI06 & 1 & 3 & 3 & 2 & 2 & 73,3 & 3 \\
APO10 & 1 & 2 & 3 & 2 & 3 & 70,0 & 3 \\
BAI07 & 1 & 3 & 3 & 1 & 2 & 68,3 & 3 \\
MEA02 & 1 & 2 & 3 & 3 & 1 & 65,0 & 3 \\
APO04 & 1 & 0 & 0 & 0 & 0 & 8,3 & 1 \\
\end{tableaucours}

\medskip
\textbf{Calcul de DSS04.} $0{,}25 \times 3 + 0{,}25 \times 3 + 0{,}20
\times 3 + 0{,}15 \times 3 + 0{,}15 \times 1 = 2{,}70$~; $S = 2{,}70 / 3
\times 100 = 90{,}0$, d'où une cible de niveau~4.

\medskip
\textbf{Lecture.} Les notes de la colonne « cascade » reprennent les liens
du Tableau~6.7~: DSS04, EDM03 et APO12 sont liés à EG06 ou EG02, BAI06 et
BAI07 ne le sont pas, d'où leur note de~1. Ils remontent par le risque et
les problèmes. APO04 (innovation) tombe au niveau~1~: la banque n'en fait
pas une priorité pendant le redressement, ce qui ne veut pas dire qu'elle
renonce à innover.
\end{boxscenario}

\begin{boxprocedure}
\textbf{Le dossier de conception.} Pour qu'un auditeur, ou le régulateur,
puisse comprendre pourquoi un objectif a reçu telle capacité cible, le
dossier de conception doit contenir~:
\begin{itemize}
    \item la valeur retenue pour chacun des onze facteurs, avec sa
          justification~;
    \item les poids des facteurs et la décision qui les a approuvés~;
    \item les notes et les scores, avec la méthode de calcul et sa source
          (outil officiel ou méthode interne)~;
    \item le canevas de conception, avec les conflits relevés et la manière
          dont ils ont été résolus~;
    \item la liste finale des objectifs et de leurs capacités cibles,
          approuvée par le comité stratégique des SI.
\end{itemize}
\end{boxprocedure}

% ============================================================================
\section{Arbitrer et Justifier}
\label{sec:c07-arbitrage}
% ============================================================================

\subsection{Résoudre les conflits}

Le calcul propose~; l'organe de gouvernance décide. La quatrième étape de
la démarche consiste à confronter le résultat aux contraintes et aux
intentions de l'organisation.

\begin{boxscenario}{Opération SANAGA. Le conflit sur BAI06}
\textbf{Le conflit.} Le calcul donne une cible de niveau~3 pour BAI06. Le
président du comité d'audit demande le niveau~4~: « Après ce qui s'est
passé, la gestion des changements doit être exemplaire. »

\medskip
\textbf{Les arguments.} Le niveau~4 suppose un processus piloté par des
mesures quantitatives (Chapitre~\ref{chap:cobit-capacite}). La DSI estime
qu'il faut d'abord atteindre le niveau~3, c'est-à-dire un processus
défini et appliqué partout, ce qui n'est pas le cas aujourd'hui. Viser
directement le niveau~4 disperserait l'effort.

\medskip
\textbf{La décision} (procès-verbal du comité stratégique). Cible de
niveau~3 à dix-huit mois, puis réexamen pour le niveau~4. En contrepartie,
un indicateur trimestriel (part des changements majeurs dotés d'un
retour arrière testé) est présenté au comité d'audit dès le premier
trimestre.
\end{boxscenario}

\subsection{Trois types de conflits fréquents}

\begin{tableaucours}{L{3.4cm}L{4.4cm}L{4.4cm}}{%
    \tableauHeader{Conflit} &
    \tableauHeader{Exemple} &
    \tableauHeader{Piste de résolution}
}
Contrôle contre agilité & Une jeune entreprise en DevOps soumise à des
exigences de conformité élevées & Automatiser les contrôles dans la chaîne
de déploiement plutôt que d'ajouter des validations manuelles \\
Ambition contre moyens & Dix objectifs visés au niveau~4 avec une équipe
de cinq personnes & Échelonner les cibles dans le temps, concentrer
l'effort sur trois objectifs \\
Facteurs contradictoires & Stratégie d'innovation, mais profil de risque
très défavorable & Dissocier les périmètres~: innover sur un périmètre
isolé, stabiliser le cœur \\
\end{tableaucours}
\vspace{-0.8em}
\begin{center}{\small\itshape Tableau~7.2. Conflits de conception et
pistes de résolution}\end{center}

\begin{boiteRemarque}{NOUN PAY a-t-elle besoin de COBIT~?}
Avec 45~salariés, NOUN PAY se situe en dessous du seuil de 50~salariés à
partir duquel COBIT~2019 considère une organisation comme une PME~: les
très petites organisations ne sont pas couvertes par le modèle. Faut-il
pour autant renoncer~? Non, mais il faut l'utiliser comme une liste de
contrôle plutôt que comme un système complet~: retenir cinq ou six
objectifs, viser des capacités de niveau~2 ou~3, s'appuyer sur le domaine
d'intérêt consacré aux PME pour réduire les structures, et réserver
l'effort aux risques qui la menacent vraiment (dépendance au fournisseur
de nuage, données personnelles, notation par IA). C'est l'application
directe du principe~5~: adapter le système aux besoins de l'organisation.
\end{boiteRemarque}

\begin{boxmetier}{DSI}
\textbf{Présenter un profil de conception à la direction générale}
\begin{enumerate}
    \item Commencer par les facteurs, en langage métier~: « nous sommes une
          banque dont les SI sont stratégiques, sous forte contrainte
          réglementaire, qui sort d'une panne majeure ».
    \item Présenter la liste des objectifs prioritaires et leur capacité
          cible, sur une seule page.
    \item Montrer ce qui \emph{n'est pas} retenu, et pourquoi~: c'est ce qui
          rend le profil crédible.
    \item Chiffrer l'effort de mise en œuvre de chaque cible, et proposer un
          échelonnement.
    \item Demander une décision formelle, et la date de réexamen.
\end{enumerate}
\end{boxmetier}

% ============================================================================
\separateur
% ============================================================================

\begin{boxresume}
\textbf{Ce qu'il faut retenir de ce chapitre~:}
\begin{itemize}
    \item \textbf{COBIT se conçoit, il ne s'applique pas tel quel}~: les
          40~objectifs sont équivalents dans le modèle, les facteurs de
          conception les hiérarchisent.
    \item \textbf{Trois effets}~: priorité et capacité cible des objectifs,
          variation des composantes, recours à des domaines d'intérêt.
    \item \textbf{Onze facteurs}~: stratégie, objectifs d'entreprise, profil
          de risque, problèmes liés aux SI, menaces, conformité, rôle des
          SI, approvisionnement, méthodes, adoption des technologies,
          taille.
    \item \textbf{Quatre étapes}~: comprendre le contexte, déterminer le
          périmètre initial, l'affiner, conclure en résolvant les conflits
          sur un canevas de conception.
    \item \textbf{Le calcul propose, la gouvernance décide}~: la méthode de
          score aide à fixer les capacités cibles, mais les arbitrages
          relèvent du comité stratégique et se consignent au procès-verbal.
    \item \textbf{Au Cameroun}, les obligations de disponibilité et d'audit
          de la loi n°~2010/012 et la loi n°~2024/017 sur les données
          personnelles placent les exigences de conformité au niveau
          « élevé » pour les banques et les acteurs du paiement.
\end{itemize}
\end{boxresume}

\bigskip

\begin{boxexo}[title={\ding{45}\quad Exercice~7.1. Les facteurs de KRIBILOG \niveaubase}]
Renseignez, pour KRIBILOG SA (opérateur logistique portuaire, 1~200
salariés, exploitation confiée à l'infogérant NETPORT SERVICES), la valeur
de chacun des onze facteurs de conception, sur le modèle du Tableau~7.1.
Justifiez en une phrase les facteurs 4 (problèmes liés aux SI), 7 (rôle
des SI) et 8 (approvisionnement).
\end{boxexo}

\begin{boxexo}[title={\ding{45}\quad Exercice~7.2. Le profil de l'ANPP \niveaumoyen}]
Le comité stratégique de l'ANPP retient les poids suivants~: cascade
(30~\%), problèmes liés aux SI (30~\%), exigences de conformité (20~\%),
approvisionnement (20~\%). Les notes proposées sont~:
EDM02~(3,~3,~2,~1), APO05~(3,~3,~1,~1), APO10~(1,~3,~2,~3),
BAI11~(2,~3,~1,~2), DSS04~(1,~1,~2,~1).
\begin{enumerate}[(a)]
    \item Calculez le score d'importance de chaque objectif.
    \item Déduisez-en la capacité cible selon la convention du cours.
    \item Le directeur général demande que DSS04 soit porté au niveau~3.
          Quels arguments pourriez-vous avancer pour et contre~?
\end{enumerate}
\end{boxexo}

\begin{boxexo}[title={\ding{45}\quad Exercice~7.3. Concevoir le système de NOUN PAY \niveauexpert}]
NOUN PAY développe ses applications en DevOps, héberge tout son SI chez un
fournisseur de nuage étranger et doit se conformer, depuis juin~2026, à la
loi n°~2024/017 sur les données personnelles.
\begin{enumerate}[(a)]
    \item Renseignez les onze facteurs de conception, en signalant ceux pour
          lesquels COBIT n'est pas pleinement adapté à la taille de
          l'entreprise.
    \item Identifiez au moins deux conflits sur le canevas de conception et
          proposez leur résolution.
    \item Proposez une liste de six objectifs au plus, avec leur capacité
          cible et leur échéance.
    \item Rédigez la page de synthèse que Mme~KAMGA présentera à ses
          investisseurs.
\end{enumerate}
\end{boxexo}

\bigskip

\begin{boxinfo}
\textbf{Transition.} La conception a fixé des capacités cibles. Il reste à
savoir ce qu'elles signifient, comment mesurer la capacité réelle d'un
processus, et comment combler l'écart. C'est l'objet du
Chapitre~\ref{chap:cobit-capacite}~: le modèle de capacité, la maturité,
et la démarche de mise en œuvre en sept phases.
\end{boxinfo}
